% GENERATED FILE. Edit canonical lessons, then run npm run generate:books.
% canonical-source-hash: fnv1a64:012a7fad166e305a
% canonical-lessons: MR-C543-manjuri, MR-C543-sammati, MR-C543-hokar, MR-C543-vaidh, MR-C543-olakhpatra, MR-R543-review-536-543

\chapter{Yes or No}
\label{ch:mr-yes-or-no}
\hlchaptermodality{543}

\begin{chapteropening}
\textbf{By the end of this chapter:} \emph{I can talk about approval, consent and valid papers.}

Complete the last lesson of chapter 543: I can talk about approval, consent and valid papers.
\end{chapteropening}

\section[\texorpdfstring{\mr{मंजुरी} (mañjurī)}{मंजुरी (mañjurī)}]{\mr{मंजुरी} (mañjurī) — approval; sanction}

\label{lesson:MR-C543-manjuri}



\begin{quote}
Before the new one: say the Marathi for to heed, then the Marathi for to defy.
\end{quote}

\subsection*{You'll want to know: \mr{मंजुरी}}
\textbf{\mr{मंजुरी}} — \emph{mañjurī} — \textquotedblleft{}approval; sanction\textquotedblright{}.

\begin{quote}
\mr{नवीन} \mr{रस्त्याला} \mr{सरकारची} \mr{मंजुरी} \mr{मिळाली}.
\end{quote}

\begin{quote}
\emph{navīn rastyālā sarkārcī mañjurī miḷālī.}
\end{quote}

\begin{quote}
The new road got the government's approval.
\end{quote}

Feminine.

What is approved takes -\mr{ला}: \mr{रस्त्याला} \mr{मंजुरी}.

\begin{cousinweb}
\textbf{\mr{मंजुरी}} (\emph{mañjurī}) is a loan from Persian \emph{manẓūrī}, \textquotedblleft{}approval\textquotedblright{}, usually traced to Arabic \emph{manẓūr}, \textquotedblleft{}seen, accepted\textquotedblright{}.
\end{cousinweb}

\begin{grammarlens}[title={What you've built}]
One more word for saying what happens if.
\end{grammarlens}

\subsection*{Your turn}
\begin{itemize}
\raggedright
  \item \emph{Say it:} \emph{mañjurī}
  \item \emph{Say it:} \emph{mañjurī}, once more
  \item \emph{Say it:} \emph{navīn rastyālā sarkārcī mañjurī miḷālī.} — The new road got the government's approval.
  \item \emph{Recall:} say the Marathi for to heed, then the Marathi for to defy, then say \emph{mañjurī} again
\end{itemize}

\begin{tcolorbox}[breakable,colback=teal!4,colframe=teal!35!black,title={Before you move on}]
What does \mr{मंजुरी} mean? (\textquotedblleft{}Approval; sanction\textquotedblright{}.) Say it once more.
\end{tcolorbox}
\section[\texorpdfstring{\mr{संमती} (sammatī)}{संमती (sammatī)}]{\mr{संमती} (sammatī) — consent}

\label{lesson:MR-C543-sammati}



\begin{quote}
Before the new one: say the Marathi for to defy, then the Marathi for approval.
\end{quote}

\subsection*{You'll want to know: \mr{संमती}}
\textbf{\mr{संमती}} — \emph{sammatī} — \textquotedblleft{}consent\textquotedblright{}.

\begin{quote}
\mr{पालकांच्या} \mr{संमतीशिवाय} \mr{मुलं} \mr{सहलीला} \mr{जाऊ} \mr{शकत} \mr{नाहीत}.
\end{quote}

\begin{quote}
\emph{pālkāñcyā sammatīśivāy mulaṁ sahlīlā jāū śakat nāhīt.}
\end{quote}

\begin{quote}
Children cannot go on the trip without their parents' consent.
\end{quote}

Feminine.

\begin{cousinweb}
\textbf{\mr{संमती}} (\emph{sammatī}) is a tatsama, taken from Sanskrit \emph{sammati}, \textquotedblleft{}agreement, approval\textquotedblright{} (Monier-Williams).
\end{cousinweb}

\begin{grammarlens}[title={What you've built}]
One more word for saying what happens if.
\end{grammarlens}

\subsection*{Your turn}
\begin{itemize}
\raggedright
  \item \emph{Say it:} \emph{sammatī}
  \item \emph{Say it:} \emph{sammatī}, once more
  \item \emph{Say it:} \emph{pālkāñcyā sammatīśivāy mulaṁ sahlīlā jāū śakat nāhīt.} — Children cannot go on the trip without their parents' consent.
  \item \emph{Recall:} say the Marathi for to defy, then the Marathi for approval, then say \emph{sammatī} again
\end{itemize}

\begin{tcolorbox}[breakable,colback=teal!4,colframe=teal!35!black,title={Before you move on}]
What does \mr{संमती} mean? (\textquotedblleft{}Consent\textquotedblright{}.) Say it once more.
\end{tcolorbox}
\section[\texorpdfstring{\mr{होकार} (hokār)}{होकार (hokār)}]{\mr{होकार} (hokār) — assent; a yes}

\label{lesson:MR-C543-hokar}



\begin{quote}
Before the new one: say the Marathi for approval, then the Marathi for consent.
\end{quote}

\subsection*{You'll want to know: \mr{होकार}}
\textbf{\mr{होकार}} — \emph{hokār} — \textquotedblleft{}assent; a yes\textquotedblright{}.

\begin{quote}
\mr{मीराने} \mr{लग्नाला} \mr{होकार} \mr{दिला}.
\end{quote}

\begin{quote}
\emph{mīrāne lagnālā hokār dilā.}
\end{quote}

\begin{quote}
Meera said yes to the marriage.
\end{quote}

Masculine.

\mr{होकार} \mr{देणे}, to say yes; its opposite is \mr{नकार}.

\begin{cousinweb}
\textbf{\mr{होकार}} (\emph{hokār}) is a Marathi derivation built on \mr{हो}, \textquotedblleft{}yes\textquotedblright{}, with -\mr{कार}, \textquotedblleft{}a saying\textquotedblright{}, from Sanskrit \emph{kāra} (Monier-Williams), formed like \mr{नकार}, \textquotedblleft{}a no\textquotedblright{}.
\end{cousinweb}

\begin{grammarlens}[title={What you've built}]
One more word for saying what happens if.
\end{grammarlens}

\subsection*{Your turn}
\begin{itemize}
\raggedright
  \item \emph{Say it:} \emph{hokār}
  \item \emph{Say it:} \emph{hokār}, once more
  \item \emph{Say it:} \emph{mīrāne lagnālā hokār dilā.} — Meera said yes to the marriage.
  \item \emph{Recall:} say the Marathi for approval, then the Marathi for consent, then say \emph{hokār} again
\end{itemize}

\begin{tcolorbox}[breakable,colback=teal!4,colframe=teal!35!black,title={Before you move on}]
What does \mr{होकार} mean? (\textquotedblleft{}Assent; a yes\textquotedblright{}.) Say it once more.
\end{tcolorbox}
\section[\texorpdfstring{\mr{वैध} (vaidh)}{वैध (vaidh)}]{\mr{वैध} (vaidh) — valid}

\label{lesson:MR-C543-vaidh}



\begin{quote}
Before the new one: say the Marathi for consent, then the Marathi for assent.
\end{quote}

\subsection*{You'll want to know: \mr{वैध}}
\textbf{\mr{वैध}} — \emph{vaidh} — \textquotedblleft{}valid\textquotedblright{}.

\begin{quote}
\mr{माझं} \mr{पारपत्र} \mr{पुढच्या} \mr{वर्षापर्यंत} \mr{वैध} \mr{आहे}.
\end{quote}

\begin{quote}
\emph{mājhaṁ pārpatra puḍhcyā varṣāparyant vaidh āhe.}
\end{quote}

\begin{quote}
My passport is valid until next year.
\end{quote}

\begin{cousinweb}
\textbf{\mr{वैध}} (\emph{vaidh}) is a tatsama, taken from Sanskrit \emph{vaidha}, \textquotedblleft{}prescribed by rule, legal\textquotedblright{} (Monier-Williams), built on \emph{vidhi}, \textquotedblleft{}rule\textquotedblright{}.
\end{cousinweb}

\begin{grammarlens}[title={What you've built}]
One more word for saying what happens if.
\end{grammarlens}

\subsection*{Your turn}
\begin{itemize}
\raggedright
  \item \emph{Say it:} \emph{vaidh}
  \item \emph{Say it:} \emph{vaidh}, once more
  \item \emph{Say it:} \emph{mājhaṁ pārpatra puḍhcyā varṣāparyant vaidh āhe.} — My passport is valid until next year.
  \item \emph{Recall:} say the Marathi for consent, then the Marathi for assent, then say \emph{vaidh} again
\end{itemize}

\begin{tcolorbox}[breakable,colback=teal!4,colframe=teal!35!black,title={Before you move on}]
What does \mr{वैध} mean? (\textquotedblleft{}Valid\textquotedblright{}.) Say it once more.
\end{tcolorbox}
\section[\texorpdfstring{\mr{ओळखपत्र} (oḷakhpatra)}{ओळखपत्र (oḷakhpatra)}]{\mr{ओळखपत्र} (oḷakhpatra) — an identity card}

\label{lesson:MR-C543-olakhpatra}



\begin{quote}
Before the new one: say the Marathi for assent, then the Marathi for valid.
\end{quote}

\subsection*{You'll want to know: \mr{ओळखपत्र}}
\textbf{\mr{ओळखपत्र}} — \emph{oḷakhpatra} — \textquotedblleft{}an identity card\textquotedblright{}.

\begin{quote}
\mr{प्रवेशद्वारावर} \mr{सगळ्यांनी} \mr{ओळखपत्र} \mr{दाखवलं}.
\end{quote}

\begin{quote}
\emph{praveśadvārāvar sagḷyānnī oḷakhpatra dākhavlaṁ.}
\end{quote}

\begin{quote}
Everyone showed an identity card at the entrance.
\end{quote}

Neuter.

\begin{cousinweb}
\textbf{\mr{ओळखपत्र}} (\emph{oḷakhpatra}) is a Marathi compound built from \mr{ओळख}, \textquotedblleft{}knowing someone\textquotedblright{}, and \mr{पत्र}, \textquotedblleft{}a leaf, a letter\textquotedblright{}, from Sanskrit \emph{patra} (Monier-Williams).
\end{cousinweb}

\begin{grammarlens}[title={What you've built}]
One more word for saying what happens if.
\end{grammarlens}

\subsection*{Your turn}
\begin{itemize}
\raggedright
  \item \emph{Say it:} \emph{oḷakhpatra}
  \item \emph{Say it:} \emph{oḷakhpatra}, once more
  \item \emph{Say it:} \emph{praveśadvārāvar sagḷyānnī oḷakhpatra dākhavlaṁ.} — Everyone showed an identity card at the entrance.
  \item \emph{Recall:} say the Marathi for assent, then the Marathi for valid, then say \emph{oḷakhpatra} again
\end{itemize}

\begin{tcolorbox}[breakable,colback=teal!4,colframe=teal!35!black,title={Before you move on}]
What does \mr{ओळखपत्र} mean? (\textquotedblleft{}An identity card\textquotedblright{}.) Say it once more.
\end{tcolorbox}
\section[(dialogue)]{Looking back: from Music Making to Yes or No}

\label{lesson:MR-R543-review-536-543}



\begin{quote}
Say the Marathi for singing as an art and an instrumentalist.
\end{quote}

\subsection*{Your turn}
Say these aloud:

\begin{itemize}
\raggedright
  \item singing as an art, an instrumentalist, a musical note, to hum, to practise music
  \item acting, a role, a rehearsal, a costume, to portray
  \item rhythm, a dancer, ankle bells, to keep the beat, to sway
  \item a publisher, an autobiography, an exhibition, a film director, to applaud
  \item insurance, a risk, a guarantee, an instalment, to be affordable
  \item caution, a precaution, vigilance, alert, to issue a warning
  \item a rulebook, control, a stern warning, to heed, to defy
  \item approval, consent, assent, valid, an identity card
\end{itemize}

\begin{tcolorbox}[breakable,colback=teal!4,colframe=teal!35!black,title={Before you move on}]
Singing as an art? (\textbf{\mr{गायन}}, \emph{gāyan}.) An instrumentalist? (\textbf{\mr{वादक}}, \emph{vādak}.) A musical note? (\textbf{\mr{सूर}}, \emph{sūr}.) To hum? (\textbf{\mr{गुणगुणणे}}, \emph{guṇguṇṇe}.) To practise music? (\textbf{\mr{रियाज} \mr{करणे}}, \emph{riyāj karṇe}.) Acting? (\textbf{\mr{अभिनय}}, \emph{abhinay}.) A role? (\textbf{\mr{भूमिका}}, \emph{bhūmikā}.) A rehearsal? (\textbf{\mr{तालीम}}, \emph{tālīm}.) A costume? (\textbf{\mr{वेशभूषा}}, \emph{veśbhūṣā}.) To portray? (\textbf{\mr{साकारणे}}, \emph{sākārṇe}.) Rhythm? (\textbf{\mr{ताल}}, \emph{tāl}.) A dancer? (\textbf{\mr{नर्तक}}, \emph{nartak}.) Ankle bells? (\textbf{\mr{घुंगरू}}, \emph{ghuṅgrū}.) To keep the beat? (\textbf{\mr{ठेका} \mr{धरणे}}, \emph{ṭhekā dharṇe}.) To sway? (\textbf{\mr{डोलणे}}, \emph{ḍolṇe}.) A publisher? (\textbf{\mr{प्रकाशक}}, \emph{prakāśak}.) An autobiography? (\textbf{\mr{आत्मचरित्र}}, \emph{ātmacaritra}.) An exhibition? (\textbf{\mr{प्रदर्शन}}, \emph{pradarśan}.) A film director? (\textbf{\mr{दिग्दर्शक}}, \emph{digdarśak}.) To applaud? (\textbf{\mr{दाद} \mr{देणे}}, \emph{dād deṇe}.) Insurance? (\textbf{\mr{विमा}}, \emph{vimā}.) A risk? (\textbf{\mr{जोखीम}}, \emph{jokhīm}.) A guarantee? (\textbf{\mr{हमी}}, \emph{hamī}.) An instalment? (\textbf{\mr{हप्ता}}, \emph{haptā}.) To be affordable? (\textbf{\mr{परवडणे}}, \emph{parvaḍṇe}.) Caution? (\textbf{\mr{सावधगिरी}}, \emph{sāvadhgirī}.) A precaution? (\textbf{\mr{खबरदारी}}, \emph{khabardārī}.) Vigilance? (\textbf{\mr{दक्षता}}, \emph{dakṣatā}.) Alert? (\textbf{\mr{सतर्क}}, \emph{satark}.) To issue a warning? (\textbf{\mr{इशारा} \mr{देणे}}, \emph{iśārā deṇe}.) A rulebook? (\textbf{\mr{नियमावली}}, \emph{niyamāvalī}.) Control? (\textbf{\mr{नियंत्रण}}, \emph{niyantraṇ}.) A stern warning? (\textbf{\mr{ताकीद}}, \emph{tākīd}.) To heed? (\textbf{\mr{जुमानणे}}, \emph{jumānṇe}.) To defy? (\textbf{\mr{झुगारणे}}, \emph{jhugārṇe}.) Approval? (\textbf{\mr{मंजुरी}}, \emph{mañjurī}.) Consent? (\textbf{\mr{संमती}}, \emph{sammatī}.) Assent? (\textbf{\mr{होकार}}, \emph{hokār}.) Valid? (\textbf{\mr{वैध}}, \emph{vaidh}.) An identity card? (\textbf{\mr{ओळखपत्र}}, \emph{oḷakhpatra}.)
\end{tcolorbox}
